\section{Kajian Pustaka}

\subsection{Hipotesis Efisiensi Proteom}

Basan dkk.\ \cite{basan2015} menunjukkan bahwa pada \textit{E.\ coli},
fraksi proteom glikolitik meningkat secara linier dengan laju pertumbuhan,
mendukung gagasan bahwa sel mengoptimalkan alokasi protein terbatas.
Pfeiffer dkk.\ \cite{pfeiffer2001} sebelumnya telah memformulasikan
trade-off hasil-versus-laju ATP: respirasi memberikan hasil tinggi
($\gamma_R \approx 26$--$32$ mol ATP/mol glukosa) namun lambat per massa
enzim, sementara glikolisis memberikan hasil rendah ($\gamma_G = 2$)
namun cepat per protein.

\subsection{Shen dkk.\ 2024}

Shen dkk.\ \cite{shen2024} mendefinisikan efisiensi proteom sebagai
$\varepsilon = J_\text{ATP} / m_\text{enzim}$ (fluks ATP aktual per massa
enzim \textit{terukur in-vivo}). Dengan data proteomik dan $^{13}$C-MFA
pada lima sistem biologis, mereka melaporkan $\varepsilon_R > \varepsilon_G$
secara konsisten. Temuan kunci: protein glikolitik diekspresikan
\textit{berlebih} (konstitutif, banyak yang idle), sehingga efisiensi
terealisasinya rendah --- fenomena yang mereka sebut ``\textit{proteome hedging}''.

\subsection{Kukurugya dkk.\ 2024}

Kukurugya dkk.\ \cite{kukurugya2024} mendefinisikan efisiensi proteom
sebagai \textit{kapasitas}: $V \cdot \gamma$, dengan $V$ = aktivitas
spesifik maksimal (V\textsubscript{max}) per miligram protein jalur.
Model 5-parameter mereka ($\gamma_R$, $\gamma_G$, $V_R$, $V_G$, $\Phi$)
memiliki solusi analitik unik tanpa parameter bebas. Mereka melaporkan
$V_G \cdot \gamma_G > V_R \cdot \gamma_R$ untuk ragi dan sel mamalia
(glikolisis lebih cepat per protein pada kapasitas penuh).

\subsection{Wang 2025}

Wang \cite{wang2025} mendamaikan keduanya dengan model alokasi-protein
baru yang memasukkan heterogenitas sel dan kualitas nutrisi. Model ini
menunjukkan bahwa kurva efisiensi respirasi dan fermentasi
\textit{berpotongan}: pada nutrisi berkualitas tinggi respirasi menang,
pada nutrisi rendah fermentasi menang.

\subsection{Kerangka sEnz}

Liu dkk.\ \cite{senz2024} mengembangkan sEnz, kerangka analisis
sensitivitas untuk model alokasi-protein menggunakan \textit{shadow price}.
Kami menggunakan pendekatan Sobol/Morris yang lebih sederhana dan
menerapkannya secara spesifik pada perselisihan Shen--Kukurugya.

\subsection{Posisi Paper Ini}

Wang mendamaikan via model baru; kami \textbf{mengaudit model asli}
untuk mendiagnosis bahwa keduanya mungkin tidak mengukur besaran
yang sama. Pendekatan ini komplementer --- diagnosis berbeda dari
pendamaian --- dan dapat dibandingkan langsung dengan mekanisme Wang.
